\documentclass[11pt,a4paper]{article}
\usepackage[margin=24mm]{geometry}
\usepackage{amsmath,amssymb,booktabs,microtype}
\usepackage[hidelinks]{hyperref}
\newcommand{\dd}{\mathrm d}
\title{Route G2-T: Mass Metrology Is Not Signed-Mode Readout}
\author{A bounded measurement audit of the unchanged Route-G action}
\date{23 September 2026}
\begin{document}
\maketitle
\begin{abstract}
Time of flight and an independent momentum determine mass, but neither
an independent momentum sensor for neutral $X$ nor a tagged flight time
is supplied by the existing S2-only readout. We derive the joint-data
requirements, their conditioning, and a same-action alternative using
directional elastic recoil. Opposite compact-mode signs require different
information: source correlations or a genuinely signed response.
An odd amplitude alone need not distinguish signs. No coupling,
production budget or previous exclusion is changed. This closes a
mathematical measurement audit, not a hardware design or particle discovery.
\end{abstract}

\section{Which clock and which unknown?}
Keep the conditional tree spectrum
\begin{equation}
 m_j=\Lambda\sqrt{25+j^2},\qquad
 \Delta{\cal L}=-\lambda_X H^\dagger H\sum_j|x_j|^2.
\end{equation}
$X$ is electrically neutral. Work in one laboratory inertial frame,
with a freely propagating massive particle between a known production
event and its first target interaction. This is an on-shell event
description, not an unspecified coherent finite-pulse calculation.
In particle-energy units $c=1$, except in flight lengths and times,
standard relativistic kinematics gives \cite{pdgkin}
\begin{equation}
 E^2=p^2+m^2,\quad \beta=p/E,\quad
 T=t_{\rm hit}-t_0=\frac{L}{c\beta},\quad
 \Delta\tau=\frac{T}{\gamma}=\frac{Lm}{cp}.
\end{equation}
Mass alone does not set a clock rate: equal velocities have equal
$\gamma$. At a common energy or a common momentum, different masses
generally have different velocities. Those are distinct preparation
conditions; they must not be interchanged during inference.

Eliminating $\beta$ yields
\begin{equation}\label{eq:mass}
 \boxed{m^2=p^2\left[\left(\frac{cT}{L}\right)^2-1\right].}
\end{equation}
Time alone fixes $m/p$: $(m,p)\mapsto (am,ap)$ leaves $T$ unchanged.
Measuring $p$ from an assumed $m$ and inserting it back into
\eqref{eq:mass} is circular. A measured energy is an alternative
independent constraint, $m=E\sqrt{1-\beta^2}$; an assumed monochromatic
source is not an energy measurement. Unknown $t_0$ adds another
degeneracy. Two timed positions may remove $t_0$ and determine velocity,
but still do not determine $m$ or $p$ separately.

\newpage
\section{Errors, exact boxes and an illustrative scale}
Let $V$ be the covariance of $(\ln p,\ln T,\ln L)$. Differentiating
$\ln m=\ln p+\tfrac12\ln[(cT/L)^2-1]$ gives
\begin{align}
 \dd\ln m&=\dd\ln p+\gamma^2(\dd\ln T-\dd\ln L),\\
 \frac{\sigma_m^2}{m^2}&\simeq
 V_{pp}+\gamma^4(V_{TT}+V_{LL}-2V_{TL})
       +2\gamma^2(V_{pT}-V_{pL}).\label{eq:cov}
\end{align}
This is a small-error expansion, not a Gaussian assumption or a
confidence interval. In particular,
$\sigma_T^2=\sigma_{\rm hit}^2+\sigma_{t_0}^2
-2\operatorname{Cov}(t_{\rm hit},t_0)$.
Source jitter, path uncertainty and correlations cannot be omitted.
Large $\gamma^2$ makes time-to-mass inversion poorly conditioned.

A distribution-free, finite-error check is possible. For
$|\delta p|/p\leq a$, $|\delta T|\leq b$,
$|\delta L|\leq d$, monotonicity gives
\begin{align}
 m_{\min}^2&=p^2(1-a)^2
       \left[\frac{c^2(T-b)^2}{(L+d)^2}-1\right],\\
 m_{\max}^2&=p^2(1+a)^2
       \left[\frac{c^2(T+b)^2}{(L-d)^2}-1\right].\label{eq:box}
\end{align}
Here $p(1-a)>0$, $L>d$, and $c(T-b)>L+d$ are required.
Outside this domain the data can produce negative estimates of $m^2$;
do not silently clip them to a physical mass. Actual inference should
use the response likelihood including such observations.

Take only for illustration a Higgs \emph{at rest},
$E=m_h/2=62.54$ GeV, $\Lambda=5$ GeV and $L=10$ m.
It is not an LHC boost distribution or an attainable beam.
\begin{center}\small
\begin{tabular}{rrrrr}
\toprule
$j$ & $m$ [GeV] & $p$ [GeV] & $T$ [ns] & mass box [GeV]\\
\midrule
$0$ &25.000000&57.325837&36.390395&[24.873,25.127]\\
$+1$&25.495098&57.107369&36.529609&[25.370,25.620]\\
$-1$&25.495098&57.107369&36.529609&[25.370,25.620]\\
$5$ &35.355339&51.587320&40.438423&[35.254,35.457]\\
\bottomrule
\end{tabular}
\end{center}
The boxes use $a=0.001$, $b=20$ ps and $d=1$ mm, as
\emph{uncalibrated requirements}, not measured resolutions or standard
deviations. Numerical full-precision endpoints are in the report.
The $j=0$ and $|j|=1$ times differ by $139.214$ ps.
With $a,d$ fixed, their reconstructed-mass boxes first touch at
$b=48.672$ ps. This number is not half the time gap: the momenta
differ, and the momentum/path errors also enter. It is a two-point
requirement test, not a new Gaussian scan or a collider classifier.
The two signed rows remain identical even with perfect resolution.

\newpage
\section{Neutral-particle readout: a missing measurement, not missing algebra}
Ordinary curvature tracking measures momentum through electric charge;
its magnetic bending relation does not supply $p$ for the neutral $X$
of this action \cite{pdgdet}. Published charged-track resolutions cannot
be assigned to $X$. A missing-momentum measurement of an invisible pair
also gives their sum, not the individual $p$ needed in \eqref{eq:mass}.

There is an exact same-action alternative if a \emph{directional}
elastic target readout is available. Let the target of known mass $M$
start at rest, recoil energy be $E_R>0$, recoil momentum magnitude
$q=\sqrt{E_R(E_R+2M)}$, and angle to the incident trajectory be $\theta$.
Set $Q=(E_R,\mathbf q)$ and $P=(E,\mathbf p)$, using signature $+---$.
The target shell gives $Q^2=-2ME_R$. The unchanged outgoing $X$ shell
$(P-Q)^2=P^2$ then gives
\begin{equation}
 2P\cdot Q=Q^2
 \quad\Longrightarrow\quad p q\cos\theta=(E+M)E_R .
\end{equation}
A timed trajectory determines $\beta=p/E$, hence
\begin{equation}\label{eq:direction}
 D=\beta q\cos\theta-E_R,\qquad
 \boxed{E=\frac{ME_R}{D},\quad p=\beta E,\quad
 m=\frac{ME_R\sqrt{1-\beta^2}}{D}.}
\end{equation}
Require $0<\beta<1$, $D>0$ and $E-E_R>0$. These conditions and the
outgoing shell are checked explicitly. $E_R$ alone is insufficient;
with $\beta,E_R,\theta$ all observed, the inversion is algebraically
unique. Its reconstructed $p$ is correlated with the same data, not
an additional independent measurement.

The conditioning is decisive. With $\beta,E_R,M$ fixed,
\begin{equation}
 \frac{\partial\ln m}{\partial\theta}
       =\frac{\beta q\sin\theta}{D}.
\end{equation}
For the $|j|=1$ row, $M=132(0.93149410242)$ GeV and $E_R=10$ keV,
$\theta\simeq89.9625^\circ$. Allowing a local 1\% mass error from angle
alone requires about $0.896$ arcsec angular error, even before timing,
energy and incident-direction errors. Small low-energy recoils are
nearly transverse here; this particular inversion is very sensitive
to angle. This is a requirement, not a claim about a directional sensor.

The actual retained XENON1T response is an energy-to-S2 response
\cite{xenon}; it supplies no nuclear recoil vector or tagged
source-to-$X$-interaction time. Electron drift time is not $X$ flight
time. Natural xenon further has an unobserved isotope mixture: use a
mixture likelihood, not the $A=132$ illustration as an event label.
Target-at-rest elastic kinematics is the declared approximation;
initial target motion, binding and joint instrument response would
need evaluation before applying this precision budget to data.
The rejected Higgs/xenon chain is not revived by this inversion.

\newpage
\section{Equal mass, opposite sign: a different information problem}
\subsection*{Retain the source-specific no-go}
For direct tree-level scalar Higgs production followed by the diagonal
portal, $m_{+j}=m_{-j}$ and the signed source measures are equal.
Every common time, direction or momentum response preserves that
equality. If $z$ denotes the complete measured record,
\begin{equation}
 {\cal P}(z|+j)={\cal P}(z|-j),\quad
 P(+j)=P(-j)=\tfrac12
 \quad\Longrightarrow\quad P_{\rm error}^{\rm opt}=\tfrac12 .
\end{equation}
A time gate cannot create information absent from the full record.
This is the G2-H source/channel theorem, not an all-orders theorem
for the full $\Phi$-holonomy action.

\subsection*{Existing source correlations can supply a signed reference}
The undriven G2 process obeys
$\Phi_n+X_l\longrightarrow\Phi_m+X_j$ with $n+l=m+j$.
If $n,l,m$ are actually tagged, $j=n+l-m$ is determined.
This uses a fixed orientation convention and a physically calibrated
source; an absolute sign without a reference is not an observable.
At $\alpha=1/4$, ideal $\Phi$ masses
$m_{\Phi,n}^2=M_5^2+(n+\alpha)^2/R^2$ identify $n$ uniquely:
equality for distinct $n,n'$ would require $n+n'=-2\alpha$.
At $2\alpha\in\mathbb Z$, this uniqueness can fail.

For unobserved $l$ one cannot just assume its value. G2-R already
derived the stronger undriven COM inversion. With observed
$r=|j|$, $p_f$, $\Phi_m$, known $n,p_i$ and $d=m-n\ne0$,
\begin{align}
 E_f&=\sqrt{m_{\Phi,m}^2+p_f^2}+\sqrt{M_D^2+r^2/R^2+p_f^2},\\
 E_n&=\sqrt{m_{\Phi,n}^2+p_i^2},\qquad
 l^2=R^2[(E_f-E_n)^2-M_D^2-p_i^2],\\
 j&=\frac{l^2-d^2-r^2}{2d}.
\end{align}
We revalidate it, not claim a newly built source. It fails to determine
a nonzero sign in elastic $d=0$ channels. Unknown incident momenta
need a source likelihood; driven events need their work exchange
$q_{\rm pump}\Omega$ included as well. This inference uses conservation
and is not independent verification of that law.

\subsection*{An odd amplitude alone is not a sign measurement}
For coherent contributions to the \emph{same} observed channel, write
$\mathcal M_{\pm}=\mathcal M_{\rm e}\pm j\mathcal M_{\rm o}$. Then
\begin{equation}\label{eq:odd}
 |{\cal M}_{+}|^2-|{\cal M}_{-}|^2
 =4j\,\operatorname{Re}({\cal M}_{\rm e}^*{\cal M}_{\rm o}).
\end{equation}
Purely odd amplitudes, or a relative phase of $\pi/2$, remain rate-blind.
Orthogonal unobserved final states do not interfere merely because
their masses coincide. A proposed new coupling must produce a surviving
interference observable, or a signed displacement in a calibrated
oriented field. Neither is activated here; gauge completion, visible
reference, constraints and source preparation would be separate work.

\newpage
\section{Time selection changes the sample, not the source budget}
Let $\xi=(\mathbf p,t_0,L,\ldots)$ have a general, possibly correlated,
incoherent production measure $\dd F_j(\xi)$ counting particles.
Let $0\leq A_j(\xi)\leq1$ include transport, scattering and preselection,
and $R_j(z|\xi)$ be a normalized joint response for such events.
For an additional gate $0\leq G(z)\leq1$,
\begin{equation}\label{eq:measure}
 \dd N_G(z)=\sum_j\int\dd F_j(\xi)\,
       A_j(\xi)R_j(z|\xi)G(z)\,\dd z .
\end{equation}
The ideal arrival component has $t=t_0+L/[c\beta_j(p)]$.
There is no assumption that $t_0$ is independent of momentum.
Marginalizing \eqref{eq:measure} over accepted arrival times can
reweight the measured momentum spectrum although free propagation
has not changed the individual momenta. A four-point, non-Gaussian
correlated-source example verifies this in the numerical report.
This is not an actual collider distribution. Finite coherent packets
would require a separately specified production/detection amplitude.

Positivity gives the distribution-independent inequality
\begin{equation}
 N_G\leq N_{G=1}\leq N_{\rm G2-H}^{\max}
                =3.9596863\times10^{-6}.
\end{equation}
The last step retains exactly G2-H's assumptions: Run-2 on-shell Higgs
budget, same portal/coupling bounds, single-pass nonmultiplying xenon
target, and the admitted $0.7$--$50$ keV elastic nuclear-recoil component.
That bound allowed \emph{every} momentum and efficiency one.
Time cuts can improve background rejection or conditional
classification, not the total signal ceiling. Changing source,
interaction, target or channel is a new physical hypothesis.

\section{Decision and reproducibility}
The bounded audit is complete. It does not establish an actual momentum
sensor, source timestamp, directional calibration or signed preparation.
Before any new instrument forecast, supply those inputs and an event
budget that survives detection. For the rejected source, production
observables avoiding a second rare scattering remain the better lead.
If signed conversion is indispensable, first test a calibrated
$\Phi$-source correlation; only then assess an explicitly proposed
new interaction. Do not add an unmeasured mode-sign coupling silently.

\texttt{code/verify\_g2\_tof\_mode\_audit.py} runs 97 checks: mass shells,
Jacobian/correlated non-Gaussian variance, finite boxes, gate monotonicity,
directional inversion and conditioning, the old source-sign inversion,
interference counterexamples, and unchanged dependency hashes.
These are algebra/implementation checks, not empirical validation.
The action and previous result files are not modified by the verifier.

\begin{thebibliography}{9}\small\raggedright
\bibitem{pdgkin} Particle Data Group, \emph{Kinematics}, 2025 update,
\href{https://pdg.lbl.gov/2025/reviews/rpp2025-rev-kinematics.pdf}
{Secs.\ 49.1, 49.4}.
\bibitem{pdgdet} Particle Data Group, \emph{Particle Detectors at
Accelerators}, 2025 update,
\href{https://pdg.lbl.gov/2025/reviews/rpp2025-rev-particle-detectors-accel.pdf}
{Sec.\ 35.12, Eq.\ (35.66)}.
\bibitem{xenon} XENON Collaboration, \emph{Light Dark Matter Search with
Ionization Signals in XENON1T},
\href{https://arxiv.org/abs/1907.11485v2}{arXiv:1907.11485v2}.
The pinned response/provenance and source budget are recorded in G2-X/H.
\end{thebibliography}
\end{document}
